We face the prospect of creating superhuman (or otherwise very powerful) AI systems in the future where those systems hold significant power in the real world \cite{Superintelligence,russell2019human}. Building up theoretical foundations for the study and design of such systems gives us a better chance to align them with our long-term interests. In this line of work, we study agent-like systems, i.e. systems optimizing their actions to maximize a certain utility function -- the framework behind the current state-of-the-art reinforcement learning systems and one of the major proposed models for future AI systems\footnote{Other major models include e.g. comprehensive systems of services \cite{drexler2019reframing} and "Oracle AI" or "Tool AI"~\cite{OracleAI}. However, there are concerns and ongoing research into the emergence of agency in these systems~\cite{omohundro2008basic,miller2020agi}.}.

If strong AI systems with the ability to act in the real world are ever deployed\footnote{Proposals to prevent this include e.g. boxing~\cite{Superintelligence} but as e.g. \citet{Leakproofing} argues, this may be difficult or impractical.}, it is very likely that they will have some means of deliberately manipulating their own implementation, either directly or indirectly (e.g. via manipulating the human controller, influencing the development of a future AI, exploiting their own bugs or physical limitations of the hardware, etc). While the extent of those means is unknown, even weak indirect means could be extensively exploited with sufficient knowledge, compute, modelling capabilities and time.

\smallskip
\citet{omohundro2008basic} argues that every intelligent system has a fundamental drive for goal preservation, because when the future instance of the same agent strives towards the same goal, it is more likely that the goal will be achieved. Therefore, Ohomundro argues, a rational agent should never modify into an agent optimizing different goals.

\citet{everitt} examine this question formally and arrive at the same conclusion: that the agent preserves its goals in time (as long as the agent's planning algorithm anticipates the consequences of self-modifications and uses the current utility function to evaluate different futures).\footnote{\citet{everitt}'s results hold independent of the length of the time horizon or temporal discounting (by simple utility scaling).
} However, Everitt's analysis assumes that the agent is a perfect utility maximizer (i.e. it always takes the action with the greatest expected utility), and has perfect knowledge of the environment. These assumptions are probably unattainable in any complex environment.

\smallskip
To address this, we present a theoretical analysis of a self-modifying agent with imperfect optimization ability and incomplete knowledge. We model the agent in the standard cybernetic model where the agent can be bounded-rational in two different ways. Either the agent makes suboptimal decisions (is a bounded-optimization agent) or has inaccurate knowledge. We conclude that imperfect optimization can lead to exponential deterioration of alignment through self-modification, as opposed to bounded knowledge, which does not result in future misalignment. An informal summary of the results is presented below.

\smallskip
Finally, we explicitly list and discuss the underlying assumptions that motivate the theoretical problem and analysis. In addition to clearly specifying the \emph{scope of conclusions}, the explicit problem assumptions can be used as a rough axis to map the space of viable research questions in the area; see Sections~\ref{sec:assumptions} and~\ref{sec:conclusions}.

\subsection{Summary of our results} \label{sec:summary}

The result of \citet{everitt} could be loosely interpreted to imply that agents with close to perfect rationality would either prefer not to self-modify, or would self-modify and only lose a negligible target value. 

We show that when we relax the assumption of perfect rationality, their result no longer applies. The bounded-rational agent may prefer to self-modify given the option and in doing so, become less aligned and lose a significant part of the attainable value according to its original goals.

\smallskip
We use the difference between the attainable and attained expected future value at an (arbitrarily chosen) future time point as a proxy for the degree of the agent's misalignment at that time. Specifically, for a future time $t$, we consider the value attainable from time $t$ (after the agent already ran and self-modified for $t$ time units), and we estimate the loss of value $f^t$ relative to the non-modified agent in the same environment state. Note that $f^t$ is not pre-discounted by the previous $t$ steps. See Section~\ref{sec:preliminaries} for formal definitions and Section~\ref{sec:assumptions} for motivation and discussion.

\smallskip
We consider four types of deviation from perfect rationality, see Section~\ref{sec:def_bounded_rat} for formal definitions.

\begin{itemize}
\item \emph{$\epsilon$-optimizers} make suboptimal decisions.
\item \emph{$\epsilon$-misaligned agents} have inaccurate knowledge of the human utility function.
\item \emph{$\epsilon$-ignorant agents} have inaccurate knowledge of the environment.
\item \emph{$\epsilon$-impatient agents} have inaccurate knowledge of the correct temporal discount function.
\end{itemize}

Note that for the sake of simplicity, we use a very simple model of bounded rationality where the errors are simply bounded by the error parameters $\epsilon_\bullet$; this has to be taken into account when interpreting the results. However, we suspect that the asymptotic dependence of value loss on the size of errors and time would be similar for a range of natural, realistic models of bounded rationality.

\medskip
\noindent
\textbf{Informal result statements}

\medskip
\noindent
\emph{Self-modifying $\epsilon$-optimizers} may deteriorate in future alignment and performance exponentially over time, losing exponential amount of utility compared to $\epsilon$-optimizers that do not self-modify. We show upper and tight lower bounds (by a constant) on the worst-case value loss in Theorem~\ref{thm:policy_modification}. As we decrease $\gamma$ (increase discounting), the rate at which the agent's performance deteriorates increases and the possibility of self-modification becomes a more serious problem. 

Our analysis of bounded-optimization agents is a generalization of Theorem 16 from \citet{everitt} in the sense that their result can be easily recovered by a basic measure-theoretic argument.

\medskip
\noindent
\emph{Self-modifying $\epsilon_u$-misaligned, $\epsilon_\rho$-ignorant, or $\epsilon_\gamma$-impatient perfect optimizers} can only lose the same value as non-self-modifying agents with the same irrationality bounds. This also holds for any combination of the three types of bounded knowledge. We give tight upper and lower bounds (up to a constant factor) for the worst-case performance. See Section~\ref{sec:optimization_bounds} for details.

This implies that unlike bounded-optimization agents, the performance of perfect-optimization bounded-knowledge agents does not deteriorate in time. This is because bounded-knowledge agents continue to take optimal actions with respect to their almost correct knowledge and do not self-modify in a way that would worsen their performance in their view. Therefore, the possibility of self-modification seems less dangerous in the case of bounded-knowledge agents than in the case of bounded-optimization agents.

\medskip
\noindent
A \emph{self-modifying agent with any combination of the four irrationality types} may lose value exponential in the time step $t$ when the agent optimization error parameter $\epsilon_o > 0$. We again give tight (up to a constant factor) lower bounds on the worst-case performance of such agents. See Section~\ref{sec:combining} for details.

\smallskip
Our results do not imply that every such agent will actually perform this poorly but the prospect of exponential deterioration is worrying in the long-term, even if it happens at a much slower speed than suggested by our results. We focus on worst-case analysis because it tells us whether we can have formal guarantees of the agent's behaviour -- a highly desirable property for powerful real-world autonomous systems, including a prospective AGI (artificial general intelligence) or otherwise strong AIs.

\medskip
\noindent
\textbf{Overview of formal results.} Here we summarize how much value the different types of bounded-rational agents may lose via misalignment. Note that the maximal attainable discounted value is at most $\frac{1}{1-\gamma}$ and the losses should be considered relative to that, or to the maximum attainable value in concrete scenarios. Otherwise, the values for different value of $\gamma$ are incomparable.
In all cases, the worst-case lower and upper bounds are tight up to a constant.

\medskip
\noindent
\emph{$\epsilon$-optimizer agents} -- bounded optimization, after $t$ steps of possible self-modification (Theorem~\ref{thm:policy_modification})
$$f^t_\text{opt}(\epsilon, \gamma) = \min(\frac{\epsilon}{\gamma^{t-1}}, \frac{1}{1-\gamma})$$ 

\medskip
\noindent
\emph{$\epsilon$-misaligned agents} -- inaccurate utility (Theorem~\ref{thm:inaccurate_utility})
$$f_\text{util}(\epsilon, \gamma) = \frac{2\epsilon}{1-\gamma}$$

\medskip
\noindent
\emph{$\epsilon$-ignorant agents} -- inaccurate belief (Theorem~\ref{thm:inaccurate_belief})
$$f_\text{bel}(\epsilon, \gamma) = \frac{2}{1-\gamma} - \frac{2}{1-\gamma(1-\epsilon)}$$

\smallskip
\noindent
\emph{$\epsilon$-impatient agents} -- inaccurate discounting (Theorem~\ref{thm:imprecise_discount}) Here $\gamma^*$ is the correct discount factor and $\gamma$ is the agent's incorrect discount factor.
$$f_\text{disc}(\gamma, \gamma^*) \approx \frac{2 {\gamma^*}^\frac{1}{\lg \gamma}-1}{1-\gamma^*}$$
    
\section{Assumptions and rationale}\label{sec:assumptions}

Both the statement of the problem and its relevance to AI alignment rest on a set of assumptions listed below. While this list is non-exhaustive, we try to cover the main implicit and explicit choices in our framing, and the space of alternatives. This is largely in hope of eventually finding a better, more robust theoretical framework for solving agent self-modification within the context of AI alignment, but even further negative results in the space would inform our intuitions on what aspects of self-modification make the problem harder.

\smallskip
We propose consideration of various assumptions as a framework for thinking about prospective \emph{realistic agent models that admit formal guarantees}. We invite further research and generalizations in this area, one high-level goal being to map a part of the space of agent models and assumptions that do or do not permit guarantees, eventually finding agent models that do come with meaningful guarantees. Further negative results would inform our intuitions on what aspects of the problems make it harder.

\begin{enumerate}[(i)]
    \item \emph{Bounded rationality model.} In the models of $\epsilon$-bounded-rational agents defined in Section~\ref{sec:bounded_optimization}, $\epsilon$ is generally an upper bound on the size of the optimization or knowledge error. One interpretation of our results is that value drift can happen even if the error is bounded at every step. One could argue that a more realistic scenario would assume some distribution of the size of the errors, assuming larger errors less likely or less frequent; see discussion below and in Section~\ref{sec:conclusions}.
    \item \emph{Unlimited self-modification ability.} We assume the agent is able to perform any self-modification at any time. This models the worst-case scenario when compared to a limited but still perfectly controlled self-modification. However, embedded (non-dualistic) agents in complex environments may chieve almost-unlimited self-modification from a limited ability, e.g. over a longer time span; see e.g. \cite{EmbeddedAgency}. We model the agent's self-modifications as orthogonal to actions in the environment.
    \item \emph{Modification-independence.} We assume that the agent's utility function does not explicitly reward or punish self-modifications. We also assume that self-modifications do not have any direct effect on the environment. This is captured by \Cref{def:mod_indep}.
    \item \emph{No corrigibility mechanisms.} We do not consider systems that would allow human operators to correct the system's goals, knowledge or behaviour. The problem of robust strong AI corrigibility is far from solved today and this paper can be read as a further argument for substantially more research in this direction.
    \item \emph{Worst-case analysis and bound tightness.} We focus on worst-case performance guarantees in abstracted models rather than e.g. full distributional analysis, and we show that our worst-case bounds are attainable (up to constant factors) under certain agent behaviour. Note this approach may turn out as too pessimistic or even impossible in some settings (e.g. quantum physics).
    \item \emph{Bounded value attainable per time unit.} We assume the agent obtains instantaneous utility between 0 and 1 at each time step. This is not an arbitrary choice: A constant bound on instantaneous value can be normalized to this interval. Instantaneous values bounded by a function of time $U(t)<\mu^t$ can be pre-discounted when $\gamma\mu<1$, and generally lead to infinite future values otherwise, which we disallow here to avoid foundational problems.
    \item \emph{Temporal value discounting.} We assume the agent employs some form of temporal value discounting. This could be motivated by technical or algorithmic limitations, increasing uncertainty about the future, or to avoid issues with incomparable infinite values of considered futures (see \citet{Bostrom2011} for a discussion of infinite ethics). Discounting, however, contrasts with the long-termist view; see the discussion below.
    \item \emph{Exponential discounting.} Our model assumes the agent discounts future utility exponentially, a standard assumption in artificial intelligence and the only time-invariant discounting schema \cite{strotz1955myopia} leading to consistent preferences over time.
    \item \emph{Unbounded temporal horizons.} Our analysis focuses on the long-term behaviour of the agent, in particular stability and performance from the perspective of future stakeholders (sharing the original utility function). Note that our results also to some extent apply to finite-horizon but long-running systems.

    Temporal discounting contrasts with the long-termist view: Why not model non-discounted future utility directly? Noting the motivations we mention in (vii), we agree that models of future value aggregation other than discounting would be generally better suited for long-term objectives. However, this seems to be a difficult task, as such models are neither well developed nor currently used in open-ended AI algorithms (with the obvious exception of a finite time horizon, which we propose to explore in Section~\ref{sec:conclusions}).
    
    We therefore propose a direct interpretation of our results: \emph{Assuming we implement agents that are $\epsilon$-optimizers with discounting, they may become exponentially less aligned over time. This is not the case with perfect optimizers with imperfect knowledge and discounting.}
    \item \emph{Dualistic setting.} We assume a dualistic agent and allow self-modification through special actions. This allows us to formally model one aspect of embedded agency -- at least until there are sufficient theoretical foundations of embedded agency.

    Note that in the embedded (non-dualistic) agent setting, it is not formally clear -- or possibly even definable -- what constitutes a self-modification, since there is no clear conceptual boundary between the agent and the environment, as discussed by \citet{EmbeddedAgency}.
\end{enumerate}

%\subsubsection{Assumption categories and the problem space.} AAAI
\paragraph{Assumption categories and the problem space.}

Each assumption identifies a subspace of research questions we would obtain by varying the relevant choices. These subspaces vary from very technical (e.g. concrete rationality model) to foundational (e.g. finite values and dualistic agent models). Along this axis, the assumptions and choices point to different kinds of prospective problems; we briefly describe three such categories and their prospects. See Section~\ref{sec:conclusions} for concrete proposals of future work.

\medskip
\noindent
\emph{Technical choices:} A concrete model of bounded rationality, unlimited self-modification model, and modification-independence. These are likely important for short and medium time-frames, where even eventually-diverging guarantees are useful.

We believe that many models within some realistic and sufficiently strong model classes would lead to qualitatively equivalent results in long time horizons; e.g. the agent divergence would be asymptotically exponential without external corrigibility, embedded agents in sufficiently complex environments would be able to self-modify arbitrarily over a long time (see discussion above) etc. However, these intuitions call for further verification.

\medskip
\noindent
\emph{Problem components:} No corrigibility mechanisms, unbounded time horizon, time-invariant temporal discounting, focus on the worst-case guarantees. For those, there are interesting alternatives that may yield more optimistic results. In particular, it would be valuable to explore formal models of corrigibility, perform a full probabilistic analysis of agent development, and develop long-term non-discounted finite-time settings.

\medskip
\noindent
\emph{Foundational assumptions:} Dualistic agent model and finite value of the future. Those are a standard in the area, but alternative settings may open up important and fruitful model classes and technical choices that capture currently pre-paradigmatic aspects (e.g. theory of embedded agency and non-dualistic agents).

%\subsection{Paper structure} AAAI
%
%Sections~\ref{sec:preliminaries} and~\ref{sec:def_bounded_rat} formally define the setting and different types of bounded-rational agents.
%In Sections~\ref{sec:policy_mod}, \ref{sec:optimization_bounds}, and \ref{sec:combining} follows the exposition of our results. For the reader's convenience, we have deferred the proofs together with explanation of our techniques to Sections~\ref{sec:policy_mod_techn}, \ref{sec:optimization_bounds_techn}, and \ref{sec:combining_techn}. We highlight several promising research directions in \Cref{sec:conclusions}.
